% GENERATED FILE. Edit canonical lessons, then run npm run generate:books.
% canonical-source-hash: fnv1a64:41e2e5c231eaea3d
% canonical-lessons: BN-C528-songgoti, BN-C528-phikir, BN-C528-samorthyo, BN-C528-kulono, BN-C528-jotano

\chapter{Ways and Means}
\label{ch:bn-ways-and-means}
\hlchaptermodality{528}

\begin{chapteropening}
\textbf{By the end of this chapter:} \emph{I can say whether we can afford something and find a way.}

Complete the last lesson of chapter 528: I can say whether we can afford something and find a way.
\end{chapteropening}

\section[\texorpdfstring{\bn{সঙ্গতি} (sônggôti)}{সঙ্গতি (sônggôti)}]{\bn{সঙ্গতি} (sônggôti) — financial means}

\label{lesson:BN-C528-songgoti}



\begin{quote}
Before the new one: say the Bengali for a rehearsal, then the Bengali for applause.
\end{quote}

\subsection*{You'll want to know: \bn{সঙ্গতি}}
\textbf{\bn{সঙ্গতি}} — \emph{sônggôti} — \textquotedblleft{}financial means\textquotedblright{}.

\begin{quote}
\bn{যদি} \bn{সঙ্গতি} \bn{থাকত}, \bn{তাহলে} \bn{আমরা} \bn{একটা} \bn{বড়} \bn{বাড়ি} \bn{কিনতাম।}
\end{quote}

\begin{quote}
\emph{Jôdi sônggôti thāktô, tāhôle āmrā ekṭā bôṛo bāṛi kintām.}
\end{quote}

\begin{quote}
If we had the means, we would buy a big house.
\end{quote}

\begin{cousinweb}
\textbf{\bn{সঙ্গতি}} (\emph{sônggôti}) is Sanskrit \emph{saṃgati}, \textquotedblleft{}coming together, meeting\textquotedblright{} (Monier-Williams); the sense \textquotedblleft{}what one can afford\textquotedblright{} is Bengali's own.
\end{cousinweb}

\begin{grammarlens}[title={What you've built}]
One more word for saying what happens if.
\end{grammarlens}

\subsection*{Your turn}
\begin{itemize}
\raggedright
  \item \emph{Say it:} \emph{sônggôti}
  \item \emph{Say it:} \emph{sônggôti}, once more
  \item \emph{Say it:} \emph{Jôdi sônggôti thāktô, tāhôle āmrā ekṭā bôṛo bāṛi kintām.} — If we had the means, we would buy a big house.
  \item \emph{Recall:} say the Bengali for a rehearsal, then the Bengali for applause, then say \emph{sônggôti} again
\end{itemize}

\begin{tcolorbox}[breakable,colback=teal!4,colframe=teal!35!black,title={Before you move on}]
What does \bn{সঙ্গতি} mean? (\textquotedblleft{}Financial means\textquotedblright{}.) Say it once more.
\end{tcolorbox}
\section[\texorpdfstring{\bn{ফিকির} (phikir)}{ফিকির (phikir)}]{\bn{ফিকির} (phikir) — a way out; a scheme to manage}

\label{lesson:BN-C528-phikir}



\begin{quote}
Before the new one: say the Bengali for applause, then the Bengali for financial means.
\end{quote}

\subsection*{You'll want to know: \bn{ফিকির}}
\textbf{\bn{ফিকির}} — \emph{phikir} — \textquotedblleft{}a way out; a scheme to manage\textquotedblright{}.

\begin{quote}
\bn{একটা} \bn{ফিকির} \bn{পেলে} \bn{আমরা} \bn{ঝামেলাটা} \bn{এড়াতে} \bn{পারব।}
\end{quote}

\begin{quote}
\emph{Ekṭā phikir pele āmrā jhāmelāṭā eṛāte pārbô.}
\end{quote}

\begin{quote}
If we find a way out, we'll be able to get around the trouble.
\end{quote}

\begin{cousinweb}
\textbf{\bn{ফিকির}} (\emph{phikir}) is usually traced to Arabic \emph{fikr}, \textquotedblleft{}thought\textquotedblright{}, through Persian: a clever thought that gets you out of a fix.
\end{cousinweb}

\begin{grammarlens}[title={What you've built}]
One more word for saying what happens if.
\end{grammarlens}

\subsection*{Your turn}
\begin{itemize}
\raggedright
  \item \emph{Say it:} \emph{phikir}
  \item \emph{Say it:} \emph{phikir}, once more
  \item \emph{Say it:} \emph{Ekṭā phikir pele āmrā jhāmelāṭā eṛāte pārbô.} — If we find a way out, we'll be able to get around the trouble.
  \item \emph{Recall:} say the Bengali for applause, then the Bengali for financial means, then say \emph{phikir} again
\end{itemize}

\begin{tcolorbox}[breakable,colback=teal!4,colframe=teal!35!black,title={Before you move on}]
What does \bn{ফিকির} mean? (\textquotedblleft{}A way out; a scheme to manage\textquotedblright{}.) Say it once more.
\end{tcolorbox}
\section[\texorpdfstring{\bn{সামর্থ্য} (sāmôrthyô)}{সামর্থ্য (sāmôrthyô)}]{\bn{সামর্থ্য} (sāmôrthyô) — ability; the capacity (to do something)}

\label{lesson:BN-C528-samorthyo}



\begin{quote}
Before the new one: say the Bengali for financial means, then the Bengali for a way out.
\end{quote}

\subsection*{You'll want to know: \bn{সামর্থ্য}}
\textbf{\bn{সামর্থ্য}} — \emph{sāmôrthyô} — \textquotedblleft{}ability; the capacity (to do something)\textquotedblright{}.

\begin{quote}
\bn{এই} \bn{কঠিন} \bn{কাজ} \bn{করার} \bn{সামর্থ্য} \bn{তার} \bn{আছে।}
\end{quote}

\begin{quote}
\emph{Ei kôṭhin kāj kôrār sāmôrthyô tār āchhe.}
\end{quote}

\begin{quote}
She has the ability to do this hard work.
\end{quote}

\begin{cousinweb}
\textbf{\bn{সামর্থ্য}} (\emph{sāmôrthyô}) is a tatsama, Sanskrit \emph{sāmarthya}, which Monier-Williams gives with a range of senses up to \textquotedblleft{}ability, power\textquotedblright{}.
\end{cousinweb}

\begin{grammarlens}[title={What you've built}]
One more word for saying what happens if.
\end{grammarlens}

\subsection*{Your turn}
\begin{itemize}
\raggedright
  \item \emph{Say it:} \emph{sāmôrthyô}
  \item \emph{Say it:} \emph{sāmôrthyô}, once more
  \item \emph{Say it:} \emph{Ei kôṭhin kāj kôrār sāmôrthyô tār āchhe.} — She has the ability to do this hard work.
  \item \emph{Recall:} say the Bengali for financial means, then the Bengali for a way out, then say \emph{sāmôrthyô} again
\end{itemize}

\begin{tcolorbox}[breakable,colback=teal!4,colframe=teal!35!black,title={Before you move on}]
What does \bn{সামর্থ্য} mean? (\textquotedblleft{}Ability; the capacity (to do something)\textquotedblright{}.) Say it once more.
\end{tcolorbox}
\section[\texorpdfstring{\bn{কুলোনো} (kulono)}{কুলোনো (kulono)}]{\bn{কুলোনো} (kulono) — to be enough; to suffice}

\label{lesson:BN-C528-kulono}



\begin{quote}
Before the new one: say the Bengali for a way out, then the Bengali for capacity.
\end{quote}

\subsection*{You'll want to know: \bn{কুলোনো}}
\textbf{\bn{কুলোনো}} — \emph{kulono} — \textquotedblleft{}to be enough; to suffice\textquotedblright{}.

\begin{quote}
\bn{এই} \bn{টাকায়} \bn{মাস} \bn{না} \bn{কুলোলে} \bn{আমি} \bn{ধার} \bn{নেব।}
\end{quote}

\begin{quote}
\emph{Ei ṭākāy mās nā kulole āmi dhār nebô.}
\end{quote}

\begin{quote}
If this money doesn't last the month, I'll borrow.
\end{quote}

Past (he/she …): \bn{সেই} \bn{টাকায়} \bn{তার} \bn{কুলোল} \bn{না}, \emph{sei ṭākāy tār kulolô nā}.

\begin{cousinweb}
\textbf{\bn{কুলোনো}} (\emph{kulono}) is said of money, time or food that will or will not stretch far enough; where the word comes from is not settled.
\end{cousinweb}

\begin{grammarlens}[title={What you've built}]
One more word for saying what happens if.
\end{grammarlens}

\subsection*{Your turn}
\begin{itemize}
\raggedright
  \item \emph{Say it:} \emph{kulono}
  \item \emph{Say it:} \emph{kulono}, once more
  \item \emph{Say it:} \emph{Ei ṭākāy mās nā kulole āmi dhār nebô.} — If this money doesn't last the month, I'll borrow.
  \item \emph{Recall:} say the Bengali for a way out, then the Bengali for capacity, then say \emph{kulono} again
\end{itemize}

\begin{tcolorbox}[breakable,colback=teal!4,colframe=teal!35!black,title={Before you move on}]
What does \bn{কুলোনো} mean? (\textquotedblleft{}To be enough; to suffice\textquotedblright{}.) Say it once more.
\end{tcolorbox}
\section[\texorpdfstring{\bn{জোটানো} (joṭāno)}{জোটানো (joṭāno)}]{\bn{জোটানো} (joṭāno) — to scrape together; to manage to find}

\label{lesson:BN-C528-jotano}



\begin{quote}
Before the new one: say the Bengali for capacity, then the Bengali for to suffice.
\end{quote}

\subsection*{You'll want to know: \bn{জোটানো}}
\textbf{\bn{জোটানো}} — \emph{joṭāno} — \textquotedblleft{}to scrape together; to manage to find\textquotedblright{}.

\begin{quote}
\bn{টাকা} \bn{জোটাতে} \bn{পারলে} \bn{মজুরটা} \bn{ছেলেকে} \bn{স্কুলে} \bn{পাঠাবে।}
\end{quote}

\begin{quote}
\emph{Ṭākā joṭāte pārle mojurṭā chheleke skule pāṭhābe.}
\end{quote}

\begin{quote}
If the labourer can scrape the money together, he will send his son to school.
\end{quote}

Past (he/she …): \bn{তিনি} \bn{টিকিটের} \bn{টাকা} \bn{জোটালেন}, \emph{tini ṭikiṭer ṭākā joṭālen}.

\begin{cousinweb}
\textbf{\bn{জোটানো}} (\emph{joṭāno}) is the causative of a verb CDIAL lists as Bengali \emph{juṭā}, \textquotedblleft{}to come together, to come to hand\textquotedblright{}, under a reconstructed \emph{yuṭati}, \textquotedblleft{}is joined\textquotedblright{} (CDIAL 10496): making things come to hand.
\end{cousinweb}

\begin{grammarlens}[title={What you've built}]
One more word for saying what happens if.
\end{grammarlens}

\subsection*{Your turn}
\begin{itemize}
\raggedright
  \item \emph{Say it:} \emph{joṭāno}
  \item \emph{Say it:} \emph{joṭāno}, once more
  \item \emph{Say it:} \emph{Ṭākā joṭāte pārle mojurṭā chheleke skule pāṭhābe.} — If the labourer can scrape the money together, he will send his son to school.
  \item \emph{Recall:} say the Bengali for capacity, then the Bengali for to suffice, then say \emph{joṭāno} again
\end{itemize}

\begin{tcolorbox}[breakable,colback=teal!4,colframe=teal!35!black,title={Before you move on}]
What does \bn{জোটানো} mean? (\textquotedblleft{}To scrape together; to manage to find\textquotedblright{}.) Say it once more.
\end{tcolorbox}
